\section{Method}
\label{sec:method}
%% -----------------------------------------------------------------------

\subsection{Progress Signal}

Let $\ell_i^{(t)}$ be the per-sample cross-entropy at step $t$.
EWL maintains a per-sample EMA as a smoothed loss history.
The EMA is updated \emph{before} computing the signal, so the velocity
numerator uses pre-update history $\mu^{(t-1)}$:

{\small
\begin{equation}
\begin{aligned}
  \mu_i^{(t)} &= \alpha\,\mu_i^{(t-1)} + (1-\alpha)\,\ell_i^{(t)}, \\
  s_i^{\text{rel}} &=
  \underbrace{\frac{\mu_i^{(t-1)} - \ell_i^{(t)}}{|\mu_i^{(t-1)}| + \varepsilon}}_{\text{velocity}}
  \cdot \mu_i^{(t)},
  \qquad \varepsilon = 10^{-8}.
\end{aligned}
\label{eq:rel_progress}
\end{equation}
}

Dividing by $|\mu^{(t-1)}|$ makes the signal \textbf{scale-invariant}: a 10\%
drop from loss 3.5 and a 10\% drop from loss 0.35 yield the same velocity.
Mastered samples have $\mu \approx \ell \approx 0$, so $s^{\text{rel}} \approx 0$.
Mislabelled samples have $\ell \approx \mu \gg 0$ (loss never sustainably
declines), so again $s^{\text{rel}} \approx 0$.
Only the adaptation frontier, where $\ell < \mu$ consistently, produces a
positive signal---this is the key property that drives implicit noise
suppression.

\subsection{LoRA Gradient Proxy and Weight Computation}

To gate the signal by adapter activity, we multiply by a \textbf{LoRA gradient
proxy} (computed after the \emph{previous} backward pass to avoid circular
dependency):

{\small
\begin{equation}
  g^{(t)} = \frac{1}{|\Lambda|}\sum_{\theta \in \Lambda} \|\nabla_\theta\|_F,
  \label{eq:grad_proxy}
\end{equation}
}

where $\Lambda$ is the set of all LoRA parameter tensors (each $A$ and $B$
matrix individually) across all targeted layers.
The combined signal $s_i = s_i^{\text{rel}} \times g^{(t-1)}$ is passed to a
softmax with temperature $\tau$:

{\small
\begin{equation}
  w_i = \frac{e^{s_i / \tau}}{\sum_j e^{s_j / \tau}}, \qquad
  \mathcal{L}_{\text{EWL}} = \sum_i w_i \ell_i.
  \label{eq:weights}
\end{equation}
}

Because $z$-score normalisation before the softmax would cancel $g$ entirely
($(g{\cdot}s - \overline{gs})/\sigma_{gs} = (s - \bar s)/\sigma_s$), the proxy
acts as an \textbf{adaptive temperature}: $1/\tau_{\text{eff}} = g^{(t-1)}/\tau$.
Weights are sharp during rapid adaptation and flatten toward uniform as the
adapter converges, recovering the SFT baseline gracefully.
A $K$-step warmup keeps weights uniform while the EMA accumulates history.

\begin{algorithm}[htbp]
\caption{EWL Training Step}
\label{alg:ewl}
\footnotesize
\begin{algorithmic}[1]
\Require Batch $\{(x_i, y_i)\}$, EMA state $\{\mu_i\}$,
         lagged proxy $g_{\text{prev}}$, temperature $\tau$, warmup $K$

\State Compute per-sample losses $\{\ell_i\}$

\If{step $\le K$}
    \textbf{return} uniform weights $1/B$
    \Comment{Warmup: avoid noisy early reweighting}
\EndIf

\State $\mu_i^{\text{prev}} \gets \mu_i$

\State $\mu_i \gets \alpha\,\mu_i + (1{-}\alpha)\,\ell_i$
\Comment{EMA tracks smoothed loss (sample difficulty)}

\State $s_i \gets \dfrac{\mu_i^{\text{prev}} - \ell_i}{|\mu_i^{\text{prev}}|+\varepsilon}
       \cdot\;\mu_i \cdot\;g_{\text{prev}}$
\Comment{Progress (loss drop) $\times$ difficulty $\times$ global scaling}

\State $w_i \gets \operatorname{softmax}(s/\tau)$
\Comment{Normalize into weights}

\State Backprop $\mathcal{L} = \sum_i w_i\,\ell_i$

\State $g_{\text{prev}} \gets \tfrac{1}{|\Lambda|}
       \sum_{\theta \in \Lambda}\|\nabla_\theta\|_F$
\Comment{Update proxy for next step}

\end{algorithmic}
\end{algorithm}

\vspace{-1em}
\subsection{Theoretical Grounding}
\label{sec:theory}

\paragraph{Gradient alignment.}
By a first-order Taylor expansion, a sample's loss decreases at step $t$ if
and only if $\nabla_\theta\ell_i^\top\mathbf{g}_t > 0$---its gradient is
\emph{aligned} with the current batch update~\citep{katharopoulos2018notall}.
EWL's velocity $(\mu_i^{(t-1)} - \ell_i)/|\mu_i^{(t-1)}|$ is a temporally
smoothed estimate of this alignment without requiring explicit per-sample
gradient computation ($O(B)$ backward passes per step).
Upweighting aligned samples reduces gradient estimator variance, and the
argument holds for any first-order optimiser.

\paragraph{LoRA subspace alignment.}
LoRA constrains updates to a rank-$r$ subspace $\Delta W = BA$.
A mislabelled sample's gradient often lies \emph{outside} this subspace; the
rank-$r$ constraint prevents the model from reconciling it, so its loss
stagnates.
EWL identifies this via stagnant loss velocity---the observable signature of
subspace misalignment---without knowing the adapter geometry.
This synergy predicts that EWL's benefit should \emph{decrease} as rank $r$
increases and more gradient directions become reachable (\S\ref{sec:discussion}).
